% Label recovery is research, not preparation. Every heuristic used to recover labels
% from the archive needs its own decision record in docs/decisions/ before it is
% described here.
\chapter{Data and Label Recovery}
\label{ch:data}

\section{The Archive}
\label{sec:archive}
% source: docs/decisions/0008-floorplancad-only-training.md
The industrial archive is reserved for external evaluation of models trained on
public drawings; training on the archive is outside the current experimental
design, and the drawings
are to be converted into a representation compatible with the public training data
before predictions are evaluated.
% source: docs/decisions/0008-floorplancad-only-training.md
This separates the training source from the intended application, although the
conversion procedure and the recovery of reference labels remain unresolved.
% source: docs/decisions/0008-floorplancad-only-training.md
The proposed route to reference labels is matching against the planner's
\gls{hscad} redraw, but the decision record does not establish a completed matching
procedure or its reliability, leaving archive evaluation as a planned experiment.

\todo[inline]{Describe the archive from an approved aggregate report, including drawing
coverage, drafting-style variation, conversion failures and the provenance of the
external evaluation labels. No archive statistics are available for this section.}

\section{Label Recovery}
\label{sec:label-recovery}
% source: docs/decisions/0008-floorplancad-only-training.md
The proposed archive evaluation requires reference labels obtained through
\gls{hscad} redraw matching.
% source: docs/decisions/0008-floorplancad-only-training.md
How the original geometry is matched to the redraw, and how unmatched geometry is
handled, has not yet been fixed.
% source: docs/decisions/0008-floorplancad-only-training.md
The label-recovery procedure remains to be defined.

\todo[inline]{Add the decision record defining coordinate alignment, matching criteria,
ambiguous matches and rejection rules before describing label recovery as an
implemented method. Report recovery coverage and failure reasons through generated
artifacts from a company-side run.}

\section{Hand-Verified Evaluation Set}
\label{sec:hand-verified-set}
% source: docs/decisions/0008-floorplancad-only-training.md
The external evaluation still requires a procedure for obtaining reference labels
from industrial drawings, with redraw matching identified as the proposed route.
% source: docs/decisions/0008-floorplancad-only-training.md
The current decision does not establish a completed, independently verified
archive evaluation set.
% source: docs/decisions/0010-evidence-backed-thesis-draft.md
Its reliability therefore remains outside the findings available in this draft.

\todo[inline]{Record whether a hand-verified archive evaluation set will be created. If it
is, document drawing selection, verification instructions, disagreement handling
and independence from training and model selection. Its existence and reliability
must not be assumed.}

\section{Public Datasets}
\label{sec:public-datasets}
% source: docs/decisions/0008-floorplancad-only-training.md
% source: docs/decisions/0005-floorplancad-class-mapping.md
FloorPlanCAD is the sole training source in the current design, and its downloaded
release associates geometric entities with a
\texttt{semantic-id} attribute, whose interpretation is fixed by a release-specific
class mapping.
% source: docs/decisions/0005-floorplancad-class-mapping.md
The mapping is stored in the experiment configuration.
% source: docs/decisions/0005-floorplancad-class-mapping.md
Class names could not be taken directly from the tables distributed with other
implementations, because their attribute convention and wall identifier differ
from those in the downloaded release.
% source: docs/decisions/0005-floorplancad-class-mapping.md
The adopted mapping was instead derived from the release itself, using layer names
as evidence and rendered crops where layer names could not separate classes;
walls, doors and windows are retained as separate targets, while the remaining semantic
classes, including curtain walls, are assigned to the background class
\texttt{other}.

% source: docs/decisions/0005-floorplancad-class-mapping.md
The evidence behind the mapping is uneven: wall-layer associations support the
wall mapping, whereas the separation of doors
from windows also depends on visual inspection of highlighted geometry.
% source: docs/decisions/0005-floorplancad-class-mapping.md
Some drawings place doors and windows on the same layer, and the adopted sliding-door
assignment is explicitly identified as uncertain in the decision record.
% source: docs/decisions/0005-floorplancad-class-mapping.md
Door and window results therefore inherit this limitation of the reference labels.

% source: docs/decisions/0006-floorplancad-split-by-source-group.md
Drawing identity includes the release folder as well as the file stem, because
identical names in different folders do not necessarily identify identical content.
% source: docs/decisions/0006-floorplancad-split-by-source-group.md
Exact duplicate files are removed before dataset counts or held-out folds are
formed.
% source: docs/decisions/0006-floorplancad-split-by-source-group.md
The remaining drawings are grouped by source for the split described in
\cref{sec:splits}.

\subsection{Parsing Coverage and Entity Types}
\label{sec:public-parsing}
% source: runs/2026-10-26-floorplancad-parse/metrics.json
The completed parsing run reports the release files processed, the distinct
contents parsed, cache reuse and failures, together with the total parsed entity
count shown in \cref{tab:public-parse-summary}.
% source: runs/2026-10-26-floorplancad-parse/metrics.json
The entity-type breakdown in \cref{tab:public-entity-kinds} includes text as well
as the geometric primitives retained by the reader.
% source: docs/decisions/0009-element-metrics.md
These inventory counts describe the input representation; text is subsequently
excluded from the classification metric.

\begin{table}[htbp]
  \centering
  % source: runs/2026-10-26-floorplancad-parse/metrics.json
  \caption{Public-release parsing summary, computed over distinct content where
  indicated. Cache-served files and distinct contents are separate reported
  quantities.}
  \label{tab:public-parse-summary}
  \input{tables/public_parse_summary.tex}
\end{table}

\begin{table}[htbp]
  \centering
  % source: runs/2026-10-26-floorplancad-parse/metrics.json
  \caption{Parsed entities by geometric kind, with text reported separately.
  Counts refer to distinct contents in the parsing run.}
  \label{tab:public-entity-kinds}
  \input{tables/public_entity_kinds.tex}
\end{table}

% source: runs/2026-10-26-floorplancad-parse/metrics.json
The recorded run reports no failed files, so parsing failures do not reduce the
public drawing inventory in this snapshot.
% source: runs/2026-10-26-floorplancad-parse/metrics.json
The count of entities with semantic identifiers is smaller than the total entity
count, making the distinction between missing labels and mapped target classes
relevant to the scoring convention.
% source: docs/decisions/0009-element-metrics.md
A missing identifier on a scored geometric entity maps to \texttt{other}, rather
than causing that entity to be omitted from evaluation.

\subsection{Release-Specific Class Interpretation}
\label{sec:public-class-interpretation}
% source: docs/decisions/0005-floorplancad-class-mapping.md
The adopted association between semantic identifiers and output classes is shown
in \cref{tab:public-class-mapping}, including the residual class used for identifiers
outside the target lists.
% source: runs/2026-10-26-floorplancad-class-audit/metrics.json
\Cref{tab:public-wall-layer-evidence} reports the wall identifier's association
with wall layers in the class audit, which covers release files before the
deduplication used for the experimental split.
% source: docs/decisions/0005-floorplancad-class-mapping.md
Layer association is evidence for interpreting an identifier, rather than an
independent verification of every assigned semantic label.

\begin{table}[htbp]
  \centering
  % source: docs/decisions/0005-floorplancad-class-mapping.md
  \caption{Release-specific semantic identifier mapping used for the evaluated
  classes. Curtain wall remains outside the wall target under the current mapping.}
  \label{tab:public-class-mapping}
  \input{tables/public_class_mapping.tex}
\end{table}

\begin{table}[htbp]
  \centering
  % source: runs/2026-10-26-floorplancad-class-audit/metrics.json
  \caption{Wall-layer association reported by the class audit over release files.
  This is label-interpretation evidence, rather than classifier precision.}
  \label{tab:public-wall-layer-evidence}
  \input{tables/public_wall_layer_evidence.tex}
\end{table}

\todo[inline]{Resolve the recorded licence question for use of FloorPlanCAD in the industry
partnership before making any claim about deployment suitability.}
